\documentclass[11pt,a4paper]{article}

% -- Packages 
\usepackage[utf8]{inputenc}  % source kept ASCII; accents via LaTeX commands
\usepackage[T1]{fontenc}
\usepackage{lmodern}
\usepackage[left=2.5cm,right=2.5cm,top=2.5cm,bottom=2.5cm]{geometry}
\usepackage{amsmath,amssymb}
\usepackage{booktabs}
\usepackage{graphicx}
\usepackage{float}
\usepackage{caption}
\usepackage[hidelinks,colorlinks=false]{hyperref}
\usepackage{setspace}
\usepackage{microtype}
\usepackage{parskip}
\usepackage{natbib}
\usepackage{fancyhdr}
\usepackage{titlesec}
\usepackage{appendix}
\usepackage{enumitem}
\usepackage{xcolor}

% -- Typography 
\setstretch{1.12}
\setlength{\parindent}{0pt}
\setlength{\parskip}{6pt}
\captionsetup{font=small,labelfont=bf}

\titleformat{\section}{\large\bfseries}{\thesection.}{0.5em}{}[{\vspace{-4pt}\rule{\textwidth}{0.5pt}\vspace{2pt}}]
\titleformat{\subsection}{\normalsize\bfseries}{\thesubsection.}{0.5em}{}
\titlespacing{\section}{0pt}{12pt}{4pt}
\titlespacing{\subsection}{0pt}{8pt}{2pt}

% -- Header / footer 
\pagestyle{fancy}
\fancyhf{}
\rhead{\small\textit{EMiF Project 2}}
\lhead{\small\textit{Risk Structure After COVID-19}}
\cfoot{\small\thepage}
\renewcommand{\headrulewidth}{0.4pt}

% -- Convenience macros 
\newcommand{\rhobar}{\bar{\rho}}
\newcommand{\sigmaEq}{\sigma^{\text{eq}}}
\newcommand{\lambdaOne}{\lambda_1}
\newcommand{\Keff}{K_{\text{eff}}}
\newcommand{\Dpost}{D^{\text{post}}}

% 
\begin{document}
% 

% -- Title page 
\begin{titlepage}
\centering
\vspace*{1.5cm}

{\LARGE\bfseries Has the Structure of Risk Changed Since COVID-19?}\\[0.6em]
{\large A walk through volatility, diversification and the way shocks travel}

\vspace{1.2cm}
\rule{0.7\textwidth}{0.5pt}
\vspace{0.8cm}

{\large\textit{Empirical Methods in Finance, Project 2}}\\[0.3em]
{\normalsize Professor Florian Ielpo}\\[0.2em]
{\normalsize University of Lausanne~$\mid$~Master of Science in Finance}

\vspace{1.2cm}

{\normalsize
\begin{tabular}{ll}
\multicolumn{2}{l}{\textbf{Group members:}}\\[4pt]
& Danny Al Fallouji\\
& Selim Grar\\
& Diona Avdija\\
\end{tabular}}

\vspace{1.2cm}
{\normalsize May 2026}

\vspace{1.5cm}

\begin{minipage}{0.9\textwidth}
\small
\textbf{Abstract.}~
We ask a simple question: after the COVID-19 shock of March 2020, did the structure of financial risk
actually change, or did markets just go through a violent but temporary scare? We use 14 daily series
from January 1990 to April 2026 and a small set of tools from the course, picked one by one as the data
tell us what to look at next. Our answer is \emph{partially}. The average level of volatility and the
average level of correlation did not move. What moved is the \emph{dynamics}: volatility shocks now hit
a bit harder and fade faster (GARCH likelihood-ratio test, $p=0.003$), and cross-asset correlations
react faster to shared shocks (DCC test, $p<0.001$). Factor concentration tightened only weakly: under
honest inference the average number of independent risk factors did not fall significantly, but its
variability did. Tail losses look heavier, yet our post-COVID sample is too short to confirm it at the
1\% tail. The static link between volatility and correlation did not change. The practical message: update
the parts of a risk model that describe \emph{how fast} risk moves, not the parts that describe its
average level.
\end{minipage}

\end{titlepage}

% -- Table of contents 
\tableofcontents
\thispagestyle{fancy}
\newpage

% 
\section{Introduction and Related Work}
% 

In a few weeks of March 2020 the S\&P~500 fell by more than a third, oil futures briefly traded below
zero, and assets that normally drift apart suddenly dropped together. It was the sharpest shock to
markets since 2008. Five years on, a fair question for a risk manager is whether that episode left a
mark on the \emph{structure} of risk, or whether markets simply passed through an extreme moment and
went back to behaving as before.

We need to say what ``structure of risk'' means before we can test it. We take it to be the shape of the
joint distribution of asset returns, described by a few features a portfolio manager actually cares
about: how big individual shocks are, how heavy the tails are, how much assets move together, how many
independent risk factors are left for diversification, and how shocks travel and persist over time. This
is deliberately a list of distinct features, because a market can keep the same \emph{average}
correlation while its tails get heavier or its factor structure gets tighter. Lumping everything into one
number would hide exactly the kind of change we are looking for.

Why does it matter in practice? Risk limits, diversification budgets and regulatory capital are all
calibrated on history. If the structure of risk shifted after 2020, a model fed with pre-2020 parameters
will give biased answers, and the bias will be silent. So the question is not academic curiosity; it
decides whether old calibrations can still be trusted.

The literature gives us both the hypotheses and the warnings. \citet{longin2001} and \citet{forbes2002}
show that correlations rise in turbulent periods, and that this is partly mechanical: when volatility is
high, measured correlation goes up almost by construction. That tells us to separate a genuine COVID
effect from this mechanical link. \citet{engle1982} and \citet{bollerslev1986} give us the ARCH/GARCH
family for volatility clustering, and \citet{glosten1993} the asymmetric version where bad news raises
volatility more than good news. For correlation that changes through time, \citet{engle2002} provides
the DCC model. \citet{embrechts1997} give the extreme-value tools for the tail. For regimes we lean on
\citet{hamilton1989}, and \citet{ang2002} add the useful point that correlations are themselves
regime-dependent. \citet{bae2003} push the idea further: diversification collapses when assets crash
together. The COVID papers came fast. \citet{baker2020} show the uncertainty shock was historically
large, but most of that early work stopped at the crisis weeks and never asked whether the change
stuck. That longer-horizon question is the gap we try to fill.

Our plan is a single thread, not a pile of tests. The preliminary diagnostics choose our tools, and we
then walk three blocks where each one answers the question the previous block leaves open. Block~A is
about the marginals: how big and how heavy are individual shocks? We answer it with realised volatility,
GARCH and extreme-value theory. Block~B asks whether diversification still works, using average
correlation, factor concentration and a DCC model. Block~C is about the dynamics, that is, how shocks
travel and persist, and there we use a regression, a small VAR and a regime model. We give away the
ending now. The change is partial, and it sits in the dynamics, not in the average levels.

% 
\section{Data}
% 

Everything comes from the provided \texttt{Data.xlsx} file: 14 daily series from January 1990 to April
2026, about 9{,}600 trading days. They cover five asset classes: five equity indices (S\&P~500,
EuroStoxx~50, Hang~Seng, MSCI~Emerging Markets, SMI), two ten-year sovereign yields (US and Germany),
two commodities (WTI oil and gold), three exchange rates (EUR/USD, USD/JPY, USD/CHF) and two US credit
indices (investment grade and high yield). We use no outside data.

For the price series we take log-returns, $r_t = \ln P_t - \ln P_{t-1}$. The two yields are rates, not
prices, so a log-return on them has no meaning; we use first differences, $\Delta y_t = y_t - y_{t-1}$.
Where a market is closed for a holiday while others trade, we forward-fill for at most two days to keep
the calendars aligned. From the five equity indices we build one equal-weight \emph{equity composite},
which is the ``market'' series we use throughout, so that the volatility, correlation and regression
results all describe the same object.

One honest caveat about the data. The equity composite mixes indices that close at different hours
(New York, Europe, Hong~Kong). This non-synchronous trading pushes measured same-day correlations down a
little, and the forward-fill adds a few zero-return days that do the same. Both effects are small at the
60-day horizon we work with, but we keep them in mind when reading the correlation results.

We split the sample at March~2020. The choice is easy to defend: the WHO declared the pandemic on
11~March, the Fed cut to zero on 15--17~March, equities had their worst week since 2008, and oil went
negative in April. But ``defensible date'' is not the same as ``clean break''. A Chow test on the
baseline correlation--volatility regression gives $F=0.78$ ($p=0.46$): there is no sharp jump at that
exact month. A CUSUM test on the same residuals rejects stability over the whole sample ($p<0.001$). Read
together, the picture is a gradual drift rather than one clean structural break, which is worth
remembering when we interpret the COVID dummy later.

The split is also very unbalanced: 347 pre-COVID months against 73 post-COVID months. The pre-COVID
window is not one regime; it averages the dot-com bust, the 2008 crisis and the euro crisis. So
``pre-COVID'' really means ``the long and varied past''. This imbalance is the single biggest limitation
of the study, and it shapes how carefully we read every comparison.

% 
\section{Methodology}
% 

This section explains the tools and what we expect from each before seeing the numbers. The logic is
that the diagnostics decide which tools are appropriate, so we start there.

\subsection{Diagnostics that choose the tools}

Three quick tests set up everything else. The Jarque--Bera test asks whether returns look normal; we
expect a clean no, with fat tails. A Ljung--Box test on \emph{squared} returns asks whether volatility
clusters; we expect yes. And the ADF and KPSS tests ask whether our monthly series are stationary, which
decides how we are allowed to run regressions. Fat tails point us to a Student-$t$ GARCH and to
extreme-value theory rather than to normal-based measures. Clustering justifies GARCH. And if the
correlation series turns out to be highly persistent, the course's own warning about spurious regression
(\citealp{phillips1986}: persistent series regressed in levels give high $R^2$ and autocorrelated
residuals that mean nothing) tells us to lean on a lagged-dependent (ARX) specification rather than a
plain level regression.

\subsection{An inference rule we use everywhere}

Our monthly risk measures are built on overlapping 60-day windows, so neighbouring months share most of
their data and are strongly autocorrelated. A standard $t$-test treats them as independent and would
report tiny, dishonest $p$-values. To avoid that trap we use a \emph{stationary block bootstrap} for
every pre/post comparison: we resample contiguous blocks of months, which keeps the autocorrelation, and
read the $p$-value from the resampled distribution. We also report each comparison against two
benchmarks, the full pre-COVID history and the recent 2015--2020 calm. That split lets us separate
``different from the long past'' from ``different from the recent quiet years''.

\subsection{Block A -- marginals}

We measure the monthly level of risk with realised volatility, the square root of the sum of squared
daily returns. For the dynamics we fit a GARCH(1,1) with Student-$t$ shocks to the equity composite,
\begin{equation}
r_t = \mu + \varepsilon_t,\quad \varepsilon_t = \sigma_t z_t,\quad
\sigma_t^2 = \omega + \alpha\,\varepsilon_{t-1}^2 + \beta\,\sigma_{t-1}^2 ,
\end{equation}
where $\alpha$ is how strongly today's shock feeds tomorrow's variance and $\alpha+\beta$ is how long
shocks persist. The GJR extension adds a term that lets negative shocks raise variance more than positive
ones, the leverage effect. For the tail we use peaks-over-threshold: we keep daily losses above the 95th
percentile and fit a Generalised Pareto Distribution, whose shape parameter $\xi$ measures tail
heaviness. We expect fat tails ($\xi>0$) and a leverage effect ($\gamma>0$) in both periods; the question
is whether they got stronger after COVID.

\subsection{Block B -- dependence}

For co-movement we compute, each month, the average pairwise correlation $\rhobar$ across all 14 series
on the trailing 60-day window. For factor concentration we run PCA on the same window and summarise it
with the effective number of factors, an entropy measure
\begin{equation}
\Keff = \exp\!\Big(-\textstyle\sum_k \tilde\delta_k \ln \tilde\delta_k\Big),\qquad
\tilde\delta_k = \delta_k/\textstyle\sum_j \delta_j ,
\end{equation}
where the $\delta_k$ are eigenvalues; $\Keff=14$ means perfectly spread risk, $\Keff=1$ a single driver.
Then, because an average correlation can hide a change in \emph{how} correlation moves, we fit the DCC
model of \citet{engle2002}. After filtering each equity series through its own GARCH, DCC lets the
correlation matrix evolve as
\begin{equation}
Q_t = (1-a-b)\,\bar Q + a\,z_{t-1}z_{t-1}' + b\,Q_{t-1},
\end{equation}
where $a$ is the speed at which correlations react to a shared shock and $b$ is how persistent the
correlation level is. EWMA (RiskMetrics, $\lambda=0.94$) is the same recursion with $a,b$ fixed, so we
show it as the textbook special case. We expect $a$ to rise if correlations now react faster to stress.

\subsection{Block C -- dynamics}

Finally we study volatility and correlation together. Our main regression asks whether the sensitivity of
correlation to volatility changed after COVID:
\begin{equation}
\rhobar_t = \alpha + \beta_1\,\sigmaEq_{t-1} + \beta_2\,(\sigmaEq_{t-1}\!\times\!\Dpost_t)
            + \gamma\,\Dpost_t + \varepsilon_t ,
\label{eq:modelF}
\end{equation}
where $\beta_2$ is the change in slope. Because $\rhobar$ is persistent, our \emph{primary} version adds
$\rhobar_{t-1}$ on the right (an ARX model), and all versions use Newey--West \citep{newey1987} standard
errors. A placebo version replaces the COVID dummy with a generic ``high-volatility'' dummy, to check we
are not just picking up the mechanical correlation--volatility link. We then estimate a small VAR on
(volatility, correlation) for Granger causality and impulse responses, and a two-state Markov-switching
model on volatility to see whether high-risk months cluster after 2020.

% 
\section{Empirical Results}
% 

The diagnostics come out as expected and set the stage. Jarque--Bera rejects normality for all 14 series,
and the Ljung--Box test finds volatility clustering in 13 of them, so Student-$t$ GARCH and extreme-value
tools are the right choices. On the monthly series, ADF and KPSS agree that realised volatility is
stationary but disagree on $\rhobar$ (ADF $p=0.13$, KPSS rejects): $\rhobar$ is borderline and very
persistent. That is exactly the case \citet{phillips1986} warns about, so in Block~C the ARX model is our
main specification and the level regression is only support.

Figure~\ref{fig:timeseries} shows the four monthly measures over the whole sample. The COVID line marks a
spike, but nothing in the picture looks like a permanent shift in level. The interesting changes, if any,
are in the dynamics, and that is where we now look.

\begin{figure}[H]
\centering
\includegraphics[width=\textwidth]{../results/fig1_risk_measures_timeseries.png}
\caption{The four monthly risk measures, January 1990 to April 2026. From top: realised equity volatility
$\sigmaEq$, average pairwise correlation $\rhobar$, the share of variance in the first principal component
$\lambdaOne$, and the effective number of factors $\Keff$. The red dashed line is March~2020.}
\label{fig:timeseries}
\end{figure}

\subsection{Block A -- how big and how heavy are the shocks?}

\textbf{The level did not move.} Table~\ref{tab:descriptive} compares the monthly measures before and
after COVID, with block-bootstrap $p$-values. Average realised volatility is essentially identical
(0.038 vs 0.038, $p=0.97$), and its spread is unchanged too. The violent level shock of March 2020 was
real but temporary; on a monthly average it leaves no footprint. If the level is flat, the natural next
question is whether the \emph{behaviour} of volatility changed, which is what GARCH measures.

\begin{table}[H]
\centering
\caption{Monthly risk measures before and after COVID, with block-bootstrap $p$-values.}
\label{tab:descriptive}
\small
\begin{tabular}{lrrrrrr}
\toprule
Series & Pre mean & Recent-pre & Post mean & Pre SD & Post SD & $p$(mean, full) \\
\midrule
$\sigmaEq$ (realised vol.)      & 0.038 & 0.032 & 0.038 & 0.022 & 0.021 & 0.966 \\
$\rhobar$ (avg.\ correlation)   & 0.081 & 0.118 & 0.090 & 0.069 & 0.056 & 0.639 \\
$\lambdaOne$ (PC1 share)        & 0.335 & 0.368 & 0.340 & 0.070 & 0.048 & 0.762 \\
$\Keff$ (effective factors)     & 7.692 & 7.041 & 7.250 & 1.158 & 0.614 & 0.146 \\
\bottomrule
\end{tabular}
\begin{flushleft}
\footnotesize
\textit{Notes:} Recent-pre is 2015--2020. $p$-values are two-sided stationary block bootstraps (blocks of
6 months) for the pre/post difference in means. A naive i.i.d.\ $t$-test would read close to zero for
$\Keff$; the bootstrap shows that is an artefact of overlapping windows. Pre-COVID: 347 months.
Post-COVID: 73 months.
\end{flushleft}
\end{table}

\textbf{The dynamics did change.} Table~\ref{tab:garch} reports the GARCH fits. A likelihood-ratio test
rejects equal parameters across the two periods ($\text{LR}=15.9$, $p=0.003$). Two things move in a
sensible direction. The immediate reaction to a shock, $\alpha$, is a little larger after COVID
(0.107 to 0.115), and total persistence $\alpha+\beta$ falls from 0.993 to 0.957. In plain terms, shocks
hit slightly harder and die out faster. We are careful here: the pre-COVID $\alpha+\beta$ sits almost at
one (a near-unit-root case), so any precise ``half-life'' number from it would be unstable, and we are
comparing 30 years against 6. We therefore read the direction, not a precise figure. The GJR results
(Table~\ref{tab:garch}, notes) confirm a leverage effect in both periods, slightly stronger after COVID
($\gamma$: 0.136 to 0.153), and a likelihood-ratio test that the GJR parameters changed is significant at
the 5\% level ($p=0.042$). Bad news has always raised volatility more than good news; after COVID it does
so a touch more.

\begin{table}[H]
\centering
\caption{GARCH(1,1)-$t$ on the equity composite (daily returns scaled by 100).}
\label{tab:garch}
\small
\begin{tabular}{lrrrrrr}
\toprule
Period & $\hat\omega$ & $\hat\alpha$ & $\hat\beta$ & $\hat\alpha+\hat\beta$ & $\hat\nu$ & Log-lik \\
\midrule
Full        & 0.012 & 0.112 & 0.878 & 0.990 & 6.01 & $-10{,}777$ \\
Pre-COVID   & 0.010 & 0.107 & 0.886 & 0.993 & 6.16 & $-8{,}915$ \\
Post-COVID  & 0.031 & 0.115 & 0.842 & 0.957 & 5.55 & $-1{,}855$ \\
\bottomrule
\end{tabular}
\begin{flushleft}
\footnotesize
\textit{Notes:} LR test of equal parameters pre/post: $\text{LR}=15.9$, $p=0.003$ ($\chi^2_4$). GJR
leverage term: $\hat\gamma_{\text{pre}}=0.136$, $\hat\gamma_{\text{post}}=0.153$ (both highly
significant); LR test that the GJR parameters differ pre/post: $p=0.042$.
\end{flushleft}
\end{table}

\textbf{The tail is probably heavier, but we cannot prove it.} Shocks also have tails, so we look there
directly. Table~\ref{tab:gpd} shows the GPD fit on equity-composite losses. The shape parameter more than
doubles (0.151 to 0.353) and the expected loss beyond the 1\% worst day (CVaR) rises from 3.81\% to
4.29\%. That is the right sign for ``heavier tails''. But honesty matters more than a headline: with only
82 post-COVID exceedances, the shape estimate is imprecise, and a likelihood-ratio test at the proper
95th-percentile threshold is \emph{not} significant ($p=0.39$). It only becomes significant if we lower
the threshold to the 90th percentile ($p=0.015$), which buys observations at the cost of mixing in
non-tail data. So our reading is: tail risk probably increased, but six post-COVID years are not enough to
confirm it at the 1\% tail. The Hill plot and Q--Q plots (Appendix~\ref{app:evt}) tell the same cautious
story.

\begin{table}[H]
\centering
\caption{Extreme-value (GPD) fit on equity-composite daily losses, 95th-percentile threshold.}
\label{tab:gpd}
\small
\begin{tabular}{lrrrrrr}
\toprule
Period & \# losses & \# exceed. & threshold $u$ & $\hat\xi$ & VaR$_{99\%}$ & CVaR$_{99\%}$ \\
\midrule
Pre-COVID  & 7{,}973 & 399 & 1.45\% & 0.151 & 2.75\% & 3.81\% \\
Post-COVID & 1{,}626 & ~82 & 1.31\% & 0.353 & 2.64\% & 4.29\% \\
\bottomrule
\end{tabular}
\begin{flushleft}
\footnotesize
\textit{Notes:} VaR is slightly lower post-COVID because the post threshold is lower; the heavier tail
shows up in the higher CVaR. LR test at a common 95th-percentile threshold: $p=0.39$ (68 post
exceedances, below the rule-of-thumb of 100). At the 90th percentile: $p=0.015$.
\end{flushleft}
\end{table}

So the marginals changed in part: the level is flat, the dynamics shifted, the tail is suggestive but
unproven. The obvious next question is whether assets now move together more, which would matter even more
for a diversified portfolio.

\subsection{Block B -- does diversification still work?}

\textbf{Average correlation is flat; concentration tightened only a little.} Back in
Table~\ref{tab:descriptive}, average correlation $\rhobar$ is statistically the same before and after
($p=0.64$). The effective number of factors $\Keff$ falls from 7.69 to 7.25, which looks like a loss of
diversification, but under the honest block bootstrap that drop is \emph{not} significant ($p=0.15$).
What \emph{is} significant is that its variability shrank ($p=0.02$): after COVID the market sits more
consistently at a slightly more concentrated structure, rather than at a clearly lower number of factors.
The static PCA agrees with this modest reading: seven principal components were needed to explain 80\% of
variance before COVID, six after (Figure~\ref{fig:scree}). This is a small narrowing, not a collapse, and
we report it as such. It also raises a sharper question: if the \emph{average} correlation did not move,
could the \emph{way} correlations move have changed?

\begin{figure}[H]
\centering
\includegraphics[width=0.8\textwidth]{../results/fig5_pca_scree.png}
\caption{Share of variance explained by each principal component, pre- vs post-COVID. Seven components
reach 80\% before COVID, six after, a modest tightening.}
\label{fig:scree}
\end{figure}

\textbf{Correlations now react faster.} This is the clearest result in the paper, and it comes from DCC.
Table~\ref{tab:dcc} shows the reaction parameter $a$ rising from 0.010 to 0.016 and the persistence
parameter $b$ falling from 0.988 to 0.958, with a likelihood-ratio test that strongly rejects equal
dynamics ($\text{LR}=57$, $p<0.001$). The reading is intuitive: after COVID, a shared shock pushes
correlations up faster, and the old level matters a little less. Figure~\ref{fig:ewma} plots the dynamic
correlation against the simple rolling average; both jump in March 2020, but the dynamic measure reacts
and then comes back more sharply. This is the change a single average could never show, and it is the
reason the average looked flat.

\begin{table}[H]
\centering
\caption{DCC parameters on the five equity indices (two-step estimation).}
\label{tab:dcc}
\small
\begin{tabular}{lrrr}
\toprule
Period & $\hat a$ (reaction) & $\hat b$ (persistence) & $\hat a+\hat b$ \\
\midrule
Full        & 0.010 & 0.987 & 0.997 \\
Pre-COVID   & 0.010 & 0.988 & 0.998 \\
Post-COVID  & 0.016 & 0.958 & 0.975 \\
\bottomrule
\end{tabular}
\begin{flushleft}
\footnotesize
\textit{Notes:} LR test of equal $(a,b)$ pre/post: $\text{LR}=57$, $p<0.001$ ($\chi^2_2$). EWMA
($\lambda=0.94$) is the same recursion with $a=0.06$, $b=0.94$ fixed; we use it as a cross-check
(Figure~\ref{fig:ewma}).
\end{flushleft}
\end{table}

\begin{figure}[H]
\centering
\includegraphics[width=0.9\textwidth]{../results/fig4_ewma_vs_rolling.png}
\caption{Average equity correlation: the 60-day rolling average (blue) and the DCC/EWMA dynamic measure
(orange). Both spike in March 2020; the dynamic measure reacts and reverts more sharply.}
\label{fig:ewma}
\end{figure}

So dependence did change, but in the \emph{speed} of correlation, not in its average level. That sets up
the last block. If correlations react faster, does the static link between volatility and correlation,
and the feedback between them, look any different?

\subsection{Block C -- how do shocks travel and persist?}

\textbf{The static volatility--correlation slope did not change.} Table~\ref{tab:models} reports the main
regression. In our preferred ARX specification, the COVID interaction $\beta_2$ is tiny and far from
significant ($p=0.79$), and a joint test on the COVID terms gives $p=0.76$. We are careful with words
here: this is ``we could not detect a change'', not ``the change is zero''. With 73 post-COVID months the
test has little power, and Block~B already showed the real change is in the \emph{speed} of correlation,
which a single slope cannot capture. The result holds across alternative breakpoints, a 40-day window, and
a first-difference version. The placebo (Model~G) is the useful check: replacing the COVID dummy with a
generic high-volatility dummy gives a large, significant \emph{negative} interaction ($\beta_2=-2.11$,
$p<0.001$). That is the familiar concave shape from \citet{longin2001}: when volatility is already high,
correlations sit near their ceiling and extra volatility adds little. It confirms that the
volatility--correlation link is mechanical, and that COVID did not bend it.

\begin{table}[H]
\centering
\caption{Correlation--volatility regression (dependent variable $\rhobar_t$), Newey--West errors.}
\label{tab:models}
\small
\begin{tabular}{lcccc}
\toprule
 & Model A & Model F & ARX-F (main) & Placebo G \\
\midrule
Constant                         & $0.034^{***}$ & $0.034^{***}$ & $0.006^{***}$ & $-0.016$ \\
$\sigmaEq_{t-1}$ ($\beta_1$)     & $1.301^{***}$ & $1.265^{***}$ & $0.112^{*}$   & $3.051^{***}$ \\
$\sigmaEq_{t-1}\!\times\!\Dpost$ ($\beta_2$) & -- & $0.216$ & $0.031$ & -- \\
$\Dpost$ ($\gamma$)              & -- & $0.001$ & $-0.003$ & -- \\
$\sigmaEq_{t-1}\!\times\!D^{\text{stress}}$ & -- & -- & -- & $-2.108^{***}$ \\
$\rhobar_{t-1}$ ($\phi$)         & -- & -- & $0.886^{***}$ & -- \\
\midrule
Adj.\ $R^2$ & 0.171 & 0.170 & 0.812 & 0.202 \\
$N$         & 420   & 420   & 420   & 420 \\
\bottomrule
\end{tabular}
\begin{flushleft}
\footnotesize
$^{***}p<0.01$, $^{**}p<0.05$, $^{*}p<0.10$. ARX-F is the main specification (it adds $\rhobar_{t-1}$ and
guards against spurious regression). Joint Wald test on the COVID terms in ARX-F: $p=0.76$. The full
A--H battery is in Appendix~\ref{app:models}.
\end{flushleft}
\end{table}

\textbf{The feedback loop is intact.} A VAR with four lags (chosen by Hannan--Quinn) shows two-way Granger
causality: volatility helps predict correlation ($F=4.89$, $p<0.001$) and correlation helps predict
volatility ($F=4.14$, $p=0.003$). The impulse response (Figure~\ref{fig:irf}) says a volatility shock
raises correlation for a few months and fades by about month six, and a variance decomposition attributes
14\% of correlation's forecast error to volatility at one month, peaking near 22\% at three months. The
pre-COVID Granger link is clearly significant ($p=0.006$); the post-COVID one is only borderline
($p=0.064$), but with 73 months that is what limited power looks like, so we do not lean on a pre/post
difference here. The transmission channel between volatility and correlation looks the same after COVID as
before.

\begin{figure}[H]
\centering
\includegraphics[width=0.8\textwidth]{../results/fig9_irf.png}
\caption{Response of average correlation to a one-standard-deviation volatility shock (VAR, Cholesky
ordering with volatility first). The effect peaks early and fades within about six months. Shaded band:
90\% bootstrap interval.}
\label{fig:irf}
\end{figure}

\textbf{Regimes: a weak and benchmark-dependent signal.} The two-state Markov-switching model on monthly
volatility finds a calm regime (mean 0.029) and a stressed one (mean 0.065). Whether post-COVID months sit
more often in the stressed regime depends entirely on the benchmark. Measured against the full 30-year
past, which includes 2008, the post-COVID period spends \emph{less} time in stress (19.6\% vs 26.3\%).
Measured against the recent 2015--2020 calm it spends slightly \emph{more} (19.6\% vs 18.1\%), a gap of
about 1.4 percentage points. We will not dress up a 1.4-point, benchmark-dependent difference as a structural change. On its
own, the regime evidence is too weak to support a clear conclusion (Figure~\ref{fig:ms},
Appendix~\ref{app:dynamics}).

\subsection{Putting it together}

Table~\ref{tab:synth} collects the verdict, one row per feature, with the honest confidence level. Two
features clearly moved (volatility dynamics, correlation reaction speed), two clearly did not (volatility
level, the static slope), and three are weak or unproven (tail risk, factor concentration, regimes). The
pattern is consistent: what changed is \emph{how fast} risk behaves, not its average level.

\begin{table}[H]
\centering
\caption{Did it change? A summary across the features of the risk structure.}
\label{tab:synth}
\small
\begin{tabular}{llll}
\toprule
Feature & Changed? & Main evidence & Confidence \\
\midrule
Volatility level        & No                    & realised vol mean flat ($p=0.97$)        & High \\
Volatility dynamics     & Yes                   & GARCH LR break; $\alpha\!\uparrow$, $\alpha+\beta\!\downarrow$ & Medium \\
Tail risk               & Probably              & $\xi$, CVaR up; few exceedances          & Low \\
Factor concentration    & Partly (variance)     & $\Keff$ mean flat ($p=0.15$); variance down & Low \\
Correlation dynamics    & Yes                   & DCC $a\!\uparrow$, $b\!\downarrow$; LR break & Medium-high \\
Vol.--corr.\ static slope & Not detected        & $\beta_2$ insignificant everywhere       & Limited power \\
Regime frequency        & Unclear               & benchmark-dependent sign                 & Low \\
\bottomrule
\end{tabular}
\end{table}

% 
\section{Discussion}
% 

What does this mean for someone actually managing risk? The short version is that the parts of a risk
model worth updating are the ones describing \emph{how fast} risk moves, not the ones describing its
average level.

Start with volatility. Persistence fell, so a shock fades faster than a pre-2020 model assumes. A desk
still running old GARCH parameters keeps its risk estimates high for too long after a shock and
over-hedges well into the recovery. Correlation is the subtler case. Its average did not move, yet it now
reacts faster, so a model that updates correlation slowly arrives late to a stress episode. A conditional,
DCC-style correlation reacts in time; a long rolling window does not. One warning follows from the same
results. Because the static slope did not change, a given volatility spike does not produce a
\emph{bigger} correlation jump than before, only a \emph{faster} one, and that difference is easy to miss
in a stress test.

Diversification is the mild part of the story. The number of independent risk factors did not fall in a
way we can defend statistically; it only grew more stable. A manager need not tear up the diversification
plan, but should expect a little less day-to-day movement in how concentrated the market feels. The tail
is the open question. The signal points to heavier losses, but it is not strong enough yet to act on with
confidence, so it belongs on a watch list rather than in a rewritten capital model.

% 
\section{Conclusion}
% 

Did the structure of risk change after COVID-19? Our answer is \emph{partially, and through specific
channels}. Volatility shocks became a little sharper and shorter-lived, and cross-asset correlations now
react faster to shared shocks; both changes are statistically clear. The average level of volatility and
of correlation did not move, and the static link between them did not change. Tail risk and factor
concentration point in the ``riskier'' direction but are too weak, on six post-COVID years, to confirm,
and the regime evidence depends on the benchmark.

We are upfront about the limits. The post-COVID window is only 73 months, which caps the power of every
test; we use one equity composite rather than studying each asset class separately; and by design we use
no macro or policy variables, so we describe the change without explaining its cause. The honest way to
keep this question alive is to watch two numbers as the sample grows: GARCH persistence $\alpha+\beta$ and
the DCC reaction parameter $a$. If they drift back toward their pre-2020 values, the changes we document
were temporary; if they hold, the structure of risk really did move, and the time to recalibrate is now.

% 
\newpage
\bibliographystyle{apalike}
\begin{thebibliography}{99}

\bibitem[Ang and Bekaert(2002)]{ang2002}
Ang, A.\ and Bekaert, G.\ (2002).
International asset allocation with regime shifts.
\textit{Review of Financial Studies}, 15(4), 1137--1187.

\bibitem[Bae et al.(2003)]{bae2003}
Bae, K.H., Karolyi, G.A.\ and Stulz, R.M.\ (2003).
A new approach to measuring financial contagion.
\textit{Review of Financial Studies}, 16(3), 717--763.

\bibitem[Baker et al.(2020)]{baker2020}
Baker, S.R., Bloom, N., Davis, S.J.\ and Terry, S.J.\ (2020).
COVID-induced economic uncertainty.
\textit{NBER Working Paper}, No.~26983.

\bibitem[Bollerslev(1986)]{bollerslev1986}
Bollerslev, T.\ (1986).
Generalised autoregressive conditional heteroskedasticity.
\textit{Journal of Econometrics}, 31(3), 307--327.

\bibitem[Embrechts et al.(1997)]{embrechts1997}
Embrechts, P., Kl\"{u}ppelberg, C.\ and Mikosch, T.\ (1997).
\textit{Modelling Extremal Events for Insurance and Finance}.
Springer, Berlin.

\bibitem[Engle(1982)]{engle1982}
Engle, R.F.\ (1982).
Autoregressive conditional heteroskedasticity with estimates of the variance of UK inflation.
\textit{Econometrica}, 50(4), 987--1007.

\bibitem[Engle(2002)]{engle2002}
Engle, R.F.\ (2002).
Dynamic conditional correlation: a simple class of multivariate GARCH models.
\textit{Journal of Business \& Economic Statistics}, 20(3), 339--350.

\bibitem[Forbes and Rigobon(2002)]{forbes2002}
Forbes, K.J.\ and Rigobon, R.\ (2002).
No contagion, only interdependence: measuring stock market co-movements.
\textit{Journal of Finance}, 57(5), 2223--2261.

\bibitem[Glosten et al.(1993)]{glosten1993}
Glosten, L.R., Jagannathan, R.\ and Runkle, D.E.\ (1993).
On the relation between the expected value and the volatility of the nominal excess return on stocks.
\textit{Journal of Finance}, 48(5), 1779--1801.

\bibitem[Hamilton(1989)]{hamilton1989}
Hamilton, J.D.\ (1989).
A new approach to the economic analysis of nonstationary time series and the business cycle.
\textit{Econometrica}, 57(2), 357--384.

\bibitem[Longin and Solnik(2001)]{longin2001}
Longin, F.\ and Solnik, B.\ (2001).
Extreme correlation of international equity markets.
\textit{Journal of Finance}, 56(2), 649--676.

\bibitem[Newey and West(1987)]{newey1987}
Newey, W.K.\ and West, K.D.\ (1987).
A simple, positive semi-definite, heteroskedasticity and autocorrelation consistent covariance matrix.
\textit{Econometrica}, 55(3), 703--708.

\bibitem[Phillips(1986)]{phillips1986}
Phillips, P.C.B.\ (1986).
Understanding spurious regressions in econometrics.
\textit{Journal of Econometrics}, 33(3), 311--340.

\end{thebibliography}

% 
\newpage
\begin{appendices}
% 

\section{Stationarity tests}
\label{app:adf}

\begin{table}[H]
\centering
\caption{ADF and KPSS tests on the monthly risk measures.}
\small
\begin{tabular}{lrrcrc}
\toprule
Series & ADF stat & ADF $p$ & ADF decision & KPSS $p$ & KPSS decision \\
\midrule
$\sigmaEq$   & $-5.66$ & $<0.001$ & stationary  & $\geq 0.10$ & stationary \\
$\rhobar$    & $-2.46$ & $0.126$  & unit root?  & $\leq 0.01$ & non-stationary \\
$\lambdaOne$ & $-4.12$ & $0.001$  & stationary  & $\leq 0.01$ & non-stationary \\
$\Keff$      & $-2.99$ & $0.036$  & stationary  & $\leq 0.01$ & non-stationary \\
\bottomrule
\end{tabular}
\begin{flushleft}
\footnotesize
\textit{Notes:} The disagreement for $\rhobar$ (ADF cannot reject a unit root, KPSS rejects stationarity)
is why we use the ARX and first-difference specifications in Block~C. KPSS $p$-values are reported on a
discrete grid.
\end{flushleft}
\end{table}

\section{Full regression battery}
\label{app:models}

\begin{table}[H]
\centering
\caption{Models A--H, dependent variable $\rhobar_t$, Newey--West errors.}
\small
\setlength{\tabcolsep}{5pt}
\begin{tabular}{lrrrrrrrr}
\toprule
 & A & B & C & D & E & F & G & H \\
\midrule
$\sigmaEq_{t-1}$ & $1.301^{***}$ & $1.857^{***}$ & $1.477^{***}$ & $1.830^{***}$ & $2.753^{***}$ & $1.265^{***}$ & $3.051^{***}$ & $2.307^{***}$ \\
$\times\,\Dpost$ & -- & -- & -- & -- & -- & $0.216$ & -- & $0.149$ \\
$\Dpost$         & -- & -- & -- & -- & -- & $0.001$ & -- & $-0.020$ \\
$\times\,D^{\text{stress}}$ & -- & -- & -- & -- & -- & -- & $-2.108^{***}$ & -- \\
Block B controls & & \checkmark & & & sel. & & & \checkmark \\
Block C controls & & & \checkmark & & sel. & & & \checkmark \\
Block D controls & & & & \checkmark & sel. & & & \checkmark \\
\midrule
Adj.\ $R^2$ & 0.171 & 0.287 & 0.253 & 0.261 & 0.358 & 0.170 & 0.202 & 0.409 \\
$N$ & 420 & 420 & 420 & 420 & 420 & 420 & 420 & 420 \\
\bottomrule
\end{tabular}
\begin{flushleft}
\footnotesize
$^{***}p<0.01$. Block~B: peer volatilities (oil, gold, EM, Hang~Seng). Block~C: credit spread,
transatlantic yield spread, change in US~10Y. Block~D: FX volatilities. ``sel.'' keeps the most
significant control from each block. The COVID interaction $\beta_2$ is insignificant in every model
that contains it; the only significant interaction is the placebo stress term in~G.
\end{flushleft}
\end{table}

Additional robustness (all consistent with no detectable slope change): alternative breakpoints
20~February~2020 ($\beta_2=0.19$, $p=0.50$) and 1~April~2020 ($\beta_2=0.12$, $p=0.67$); a
first-differenced dependent variable ($\beta_2=0.01$, $p=0.95$); and a 40-day rolling window. Using
$\lambdaOne$ as the dependent variable does show a significant interaction ($\beta_2=-0.82$, $p=0.02$),
the one place a COVID slope effect appears.

\section{Sub-block correlations and PC1 loadings}
\label{app:pca}

\begin{figure}[H]
\centering
\includegraphics[width=0.85\textwidth]{../results/fig2_subblock_correlations.png}
\caption{Rolling 60-day correlation between equities and each other asset class.}
\end{figure}

\begin{figure}[H]
\centering
\includegraphics[width=\textwidth]{../results/fig6_pc1_loadings.png}
\caption{Loadings on the first principal component, pre- vs post-COVID (sign fixed so equities load
positively).}
\end{figure}

\section{Tail diagnostics}
\label{app:evt}

\begin{figure}[H]
\centering
\includegraphics[width=0.9\textwidth]{../results/fig7b_hill_estimator.png}
\caption{Hill tail-index estimator against the number of upper order statistics. The post-COVID curve
sits higher, but the short sample makes the level uncertain.}
\end{figure}

\begin{figure}[H]
\centering
\includegraphics[width=0.9\textwidth]{../results/fig7_gpd_qq.png}
\caption{Q--Q plots of exceedances against the fitted GPD, pre (left) and post (right). The points track
the 45-degree line, so the GPD fit is reasonable in both periods.}
\end{figure}

\section{Regimes and transmission}
\label{app:dynamics}

\begin{figure}[H]
\centering
\includegraphics[width=0.9\textwidth]{../results/fig8_markov_switching.png}
\caption{Top: monthly realised volatility. Bottom: smoothed probability of the high-risk regime from the
two-state Markov-switching model. The red line is March~2020.}
\label{fig:ms}
\end{figure}

\begin{figure}[H]
\centering
\includegraphics[width=0.85\textwidth]{../results/fig10_fevd.png}
\caption{Forecast-error variance decomposition for the VAR on (volatility, correlation).}
\end{figure}

\begin{figure}[H]
\centering
\includegraphics[width=\textwidth]{../results/fig11_irf_prepost.png}
\caption{Correlation's response to a volatility shock, estimated separately on the pre- and post-COVID
samples. The post-COVID estimate is noisy (73 months) and we do not read a pre/post difference into it.}
\end{figure}

\end{appendices}

\end{document}
